\subsection{ Flux-surface Averaging}

\subsubsection{Surfacic densities}

We consider the flux-surface averaging of a surfacic quantity, such as a
flux of a current, generally noted ${\bf \Gamma }\left( \psi ,\theta \right) 
$. It is defined as the averaged flux of ${\bf \Gamma }$ through the
infinitesimal poloidal surface $dS\left( \psi \right) $ 
\begin{equation}
\left\langle \Gamma \right\rangle _{S}\left( \psi \right) =\frac{
\int_{dS\left( \psi \right) }{\bf dS\cdot \Gamma }\left( \psi ,\theta
\right) }{\int_{dS\left( \psi \right) }dS}
\end{equation}

In the $\left( \psi ,\theta ,\phi \right) $ system, the differential
poloidal surface element is given by (\ref{dsphipsitheta}) 
\begin{equation}
{\bf dS}=\frac{r}{\left| \nabla \psi \right| \left| \widehat{\psi }\cdot 
\widehat{r}\right| }d\psi d\theta \widehat{\phi }
\end{equation}
so that the infinitesimal poloidal surface element $dS_{p}\left( \psi
\right) $ is 
\begin{equation}
\int_{dS_{p}\left( \psi \right) }dS=\int_{dS_{p}\left( \psi \right) }\frac{r%
}{\left| \nabla \psi \right| \left| \widehat{\psi }\cdot \widehat{r}\right| }%
d\psi d\theta =d\psi \int_{0}^{2\pi }d\theta \frac{r}{\left| \nabla \psi
\right| \left| \widehat{\psi }\cdot \widehat{r}\right| }
\end{equation}
and the flux-surface averaged flux in the toroidal direction is 
\begin{equation}
\left\langle \Gamma \right\rangle _{\phi }\left( \psi \right) =\left( \frac{
dS_{p}}{d\psi }\right) ^{-1}\int_{0}^{2\pi }d\theta \frac{r}{\left| \nabla
\psi \right| \left| \widehat{\psi }\cdot \widehat{r}\right| }\left[ \widehat{
\phi }{\bf \cdot \Gamma }\left( \psi ,\theta \right) \right]
\label{fsavsurf}
\end{equation}
with 
\begin{eqnarray}
\frac{dS_{p}\left( \psi \right) }{d\psi } &=&\int_{0}^{2\pi }d\theta \frac{r%
}{\left| \nabla \psi \right| \left| \widehat{\psi }\cdot \widehat{r}\right| }
\\
&=&\int_{0}^{2\pi }d\theta \frac{1}{\left| \widehat{\psi }\cdot \widehat{r}
\right| }\frac{r}{R}\frac{1}{B_{P}}
\end{eqnarray}
Defining 
\begin{equation}
\overline{q}\left( \psi \right) \equiv \int_{0}^{2\pi }\frac{d\theta }{2\pi }%
\frac{1}{\left| \widehat{\psi }\cdot \widehat{r}\right| }\frac{r}{R}\frac{%
B_{\min }\left( \psi \right) }{B_{P}}
\end{equation}
we get 
\begin{equation}
\frac{dS_{p}\left( \psi \right) }{d\psi }=\frac{2\pi \overline{q}\left( \psi
\right) }{B_{\min }\left( \psi \right) }
\end{equation}
and 
\begin{equation}
\left\langle \Gamma \right\rangle _{\phi }\left( \psi \right) =\frac{1}{%
\overline{q}\left( \psi \right) }\int_{0}^{2\pi }\frac{d\theta }{2\pi }\frac{%
1}{\left| \widehat{\psi }\cdot \widehat{r}\right| }\frac{r}{R}\frac{B_{\min
}\left( \psi \right) }{B_{P}}\left[ \widehat{\phi }{\bf \cdot \Gamma }\left(
\psi ,\theta \right) \right]
\end{equation}

\subsubsection{Volumic densities}

We consider the flux-surface averaging of a volumic quantity, such as a
power density, generally noted ${\bf \Phi }\left( \psi ,\theta \right) $. It
is defined as the average value of ${\bf \Gamma }$ within the infinitesimal
volume $dV\left( \psi \right) $ 
\begin{equation}
\left\langle \Gamma \right\rangle _{V}\left( \psi \right) =\frac{%
\iint_{dV\left( \psi \right) }{\bf \Phi }\left( \psi ,\theta \right) dV}{%
\iint_{dV\left( \psi \right) }dV}
\end{equation}

In the $\left( \psi ,\theta ,\phi \right) $ system, the differential volume
element is given by (\ref{dvpsitheta}) 
\begin{equation}
dV=\frac{Rr}{\left| \nabla \psi \right| \left| \widehat{\psi }\cdot \widehat{%
r}\right| }d\psi d\theta d\phi
\end{equation}
so that the infinitesimal volume element $dV\left( \psi \right) $ of a
flux-surface is 
\begin{equation}
\iint_{dV\left( \psi \right) }dV=\iint_{dV\left( \psi \right) }\frac{Rr}{%
\left| \nabla \psi \right| \left| \widehat{\psi }\cdot \widehat{r}\right| }%
d\psi d\theta d\phi =d\psi \int_{0}^{2\pi }d\theta \int_{0}^{2\pi }d\phi 
\frac{Rr}{\left| \nabla \psi \right| \left| \widehat{\psi }\cdot \widehat{r}%
\right| }
\end{equation}
and the flux-surface averaged quantity in the toroidal direction is 
\begin{equation}
\left\langle {\bf \Phi }\right\rangle _{V}\left( \psi \right) =\left( \frac{%
dV}{d\psi }\right) ^{-1}\int_{0}^{2\pi }d\theta \int_{0}^{2\pi }d\phi \frac{%
Rr}{\left| \nabla \psi \right| \left| \widehat{\psi }\cdot \widehat{r}%
\right| }{\bf \Phi }\left( \psi ,\theta \right)
\end{equation}
\begin{equation}
2\pi \left( \frac{dV}{d\psi }\right) ^{-1}\frac{2\pi R_{0}\widetilde{q}}{%
B_{\min }\left( \psi \right) }
\end{equation}
with 
\begin{eqnarray}
\frac{dV\left( \psi \right) }{d\psi } &=&\int_{0}^{2\pi }d\theta
\int_{0}^{2\pi }d\phi \frac{Rr}{\left| \nabla \psi \right| \left| \widehat{%
\psi }\cdot \widehat{r}\right| } \\
&=&\int_{0}^{2\pi }d\theta \int_{0}^{2\pi }d\phi \frac{r}{\left| \widehat{%
\psi }\cdot \widehat{r}\right| }\frac{1}{B_{P}}
\end{eqnarray}

Under the assumption of axisymmetry, we get 
\begin{eqnarray}
\frac{dV\left( \psi \right) }{d\psi } &=&4\pi ^{2}\int_{0}^{2\pi }\frac{
d\theta }{2\pi }\frac{r}{\left| \widehat{\psi }\cdot \widehat{r}\right| }%
\frac{1}{B_{P}} \\
&=&\frac{4\pi ^{2}R_{0}}{B_{\min }\left( \psi \right) }\int_{0}^{2\pi }\frac{%
d\theta }{2\pi }\frac{1}{\left| \widehat{\psi }\cdot \widehat{r}\right| }%
\frac{r}{R_{0}}\frac{B_{\min }\left( \psi \right) }{B_{P}}
\end{eqnarray}

Defining 
\begin{equation}
\widehat{q}\left( \psi \right) \equiv \int_{0}^{2\pi }\frac{d\theta }{2\pi }%
\frac{1}{\left| \widehat{\psi }\cdot \widehat{r}\right| }\frac{r}{R_{0}}%
\frac{B_{\min }\left( \psi \right) }{B_{P}}
\end{equation}
we get 
\begin{equation}
\frac{dV\left( \psi \right) }{d\psi }=\frac{4\pi ^{2}R_{0}\widehat{q}\left(
\psi \right) }{B_{\min }\left( \psi \right) }
\end{equation}
and finally 
\begin{equation}
\left\langle {\bf \Phi }\right\rangle _{V}\left( \psi \right) =\frac{1}{%
\widehat{q}\left( \psi \right) }\int_{0}^{2\pi }\frac{d\theta }{2\pi }\frac{1%
}{\left| \widehat{\psi }\cdot \widehat{r}\right| }\frac{r}{R_{0}}\frac{%
B_{\min }\left( \psi \right) }{B_{P}}{\bf \Phi }\left( \psi ,\theta \right)
\end{equation}

\subsection{Current Density}

\subsubsection{Definition}

The density of current carried by electrons is given by 
\begin{equation}
{\bf J}\left( {\bf x}\right) =q_{e}\iiint d^{3}p\;{\bf v}f\left( {\bf x},%
{\bf p}\right)
\end{equation}
so that the parallel current density is 
\begin{equation}
J_{\Vert }\left( {\bf x}\right) =q_{e}\iiint d^{3}p\;v_{\Vert }f\left( {\bf x%
},{\bf p}\right)
\end{equation}
which becomes in $\left( p,\xi ,\psi ,\theta \right) $ phase space: 
\begin{equation}
J_{\Vert }\left( \psi ,\theta \right) =2\pi q_{e}\int_{0}^{\infty
}p^{2}dp\int_{-1}^{1}d\xi \;\frac{p\xi }{\gamma m}f\left( p,\xi ,\psi
,\theta \right)
\end{equation}

\subsubsection{Flux-Surface Averaging}

We are usually interested in the flux-surface averaged current density in
the toroidal direction. It is generally given by (\ref{fsavsurf}) 
\begin{eqnarray}
\left\langle J_{\Vert }\right\rangle _{\phi }\left( \psi \right) &=&\frac{1}{%
\overline{q}\left( \psi \right) }\int_{0}^{2\pi }\frac{d\theta }{2\pi }\frac{%
1}{\left| \widehat{\psi }\cdot \widehat{r}\right| }\frac{r}{R}\frac{ B_{\min
}\left( \psi \right) }{B_{P}}J_{\Vert }\left( \psi ,\theta \right) \left[ 
\widehat{\phi }{\bf \cdot }\widehat{b}\right] \\
&=&\frac{1}{\overline{q}\left( \psi \right) }\int_{0}^{2\pi }\frac{d\theta }{
2\pi }\frac{1}{\left| \widehat{\psi }\cdot \widehat{r}\right| }\frac{r}{R}%
\frac{B_{\min }\left( \psi \right) }{B_{P}}J_{\Vert }\left( \psi ,\theta
\right) \frac{B_{T}}{B}
\end{eqnarray}
and finally, using (\ref{Psidef}) 
\begin{equation}
\left\langle J_{\Vert }\right\rangle _{\phi }\left( \psi \right) =\frac{1}{%
\overline{q}\left( \psi \right) }\int_{0}^{2\pi }\frac{d\theta }{2\pi }\frac{
1}{\left| \widehat{\psi }\cdot \widehat{r}\right| }\frac{r}{R}\frac{B_{T}}{
B_{P}}\frac{J_{\Vert }\left( \psi ,\theta \right) }{\Psi \left( \psi ,\theta
\right) }
\end{equation}

\subsubsection{Fokker-Planck Equation}

When we consider only the zero order distribution function, we have that $
f_{0}$ is constant along a field line, so that 
\begin{equation}
f_{0}\left( p,\xi ,\psi ,\theta \right) =f_{0}^{(0)}\left( p,\xi _{0},\psi
\right)
\end{equation}
where 
\begin{equation}
\xi _{0}=\sigma \sqrt{1-\frac{1}{\Psi \left( \psi ,\theta \right) }\left(
1-\xi ^{2}\right) }  \label{xi0}
\end{equation}

Consequently, we find 
\begin{eqnarray}
J_{\Vert }^{0}\left( \psi ,\theta \right) &=&2\pi q_{e}\int_{0}^{\infty
}p^{2}dp\int_{-1}^{1}d\xi \;\frac{p\xi }{\gamma m}f_{0}\left( p,\xi ,\psi
,\theta \right)  \nonumber \\
&=&2\pi q_{e}\int_{0}^{\infty }p^{2}dp\int_{-1}^{1}d\xi \;\frac{p\xi }{
\gamma m}f_{0}^{(0)}\left( p,\xi _{0},\psi \right)  \nonumber \\
&=&2\pi q_{e}\int_{0}^{\infty }p^{2}dp\int_{-1}^{1}d\xi _{0}\Psi \left( \psi
,\theta \right) H\left( \left| \xi _{0}\right| -\sqrt{1-\frac{1}{\Psi \left(
\psi ,\theta \right) }}\right) \;\frac{p\xi _{0}}{\gamma m}f_{0}^{(0)}\left(
p,\xi _{0},\psi \right)  \label{jpar}
\end{eqnarray}
where the condition 
\begin{equation}
\left| \xi _{0}\right| \geq \sqrt{1-\frac{1}{\Psi \left( \psi ,\theta
\right) }}  \label{xi0cond}
\end{equation}
results from the equation (\ref{xi0}) and means that only the particle who
reach the position $\theta $ must be considered. Note that the integrand in
the equation (\ref{jpar}) is odd in $\xi _{0}$ for trapped electrons, since $%
f_{0}^{(0)}$ is symmetric in the trapped region. As a consequence, the
contribution from trapped electrons vanishes, and (\ref{jpar}) can be
rewritten as 
\begin{equation}
J_{\Vert }^{0}\left( \psi ,\theta \right) =2\pi q_{e}\int_{0}^{\infty
}p^{2}dp\int_{-1}^{1}d\xi _{0}\Psi \left( \psi ,\theta \right) H\left(
\left| \xi _{0}\right| -\xi _{0T}\right) \;\frac{p\xi _{0}}{\gamma m}%
f_{0}^{(0)}\left( p,\xi _{0},\psi \right)
\end{equation}

Therefore, the flux-surface averaged current density 
\begin{equation}
\left\langle J_{\Vert }\right\rangle _{\phi }^{0}\left( \psi \right) =\frac{1%
}{\overline{q}\left( \psi \right) }\int_{0}^{2\pi }\frac{d\theta }{2\pi }%
\frac{1}{\left| \widehat{\psi }\cdot \widehat{r}\right| }\frac{r}{R}\frac{
B_{T}}{B_{P}}\frac{J_{\Vert }^{0}\left( \psi ,\theta \right) }{\Psi \left(
\psi ,\theta \right) }
\end{equation}
becomes 
\begin{eqnarray}
\left\langle J_{\Vert }\right\rangle _{\phi }^{0}\left( \psi \right) &=&%
\frac{2\pi q_{e}}{m}\int_{0}^{\infty }dp\;\frac{p^{3}}{\gamma }\frac{1}{%
\overline{q}\left( \psi \right) }\int_{0}^{2\pi }\frac{d\theta }{2\pi }\frac{
1}{\left| \widehat{\psi }\cdot \widehat{r}\right| }\frac{r}{R}\frac{B_{T}}{
B_{P}}\frac{1}{\Psi \left( \psi ,\theta \right) } \\
&&\int_{-1}^{1}d\xi _{0}\Psi \left( \psi ,\theta \right) H\left( \left| \xi
_{0}\right| -\xi _{0T}\right) \;\xi _{0}f_{0}^{(0)}\left( p,\xi _{0},\psi
\right)
\end{eqnarray}

The integrals over $\theta $ and $\xi _{0}$ can be permuted 
\begin{eqnarray}
\left\langle J_{\Vert }\right\rangle _{\phi }^{0}\left( \psi \right) &=&%
\frac{2\pi q_{e}}{m}\int_{0}^{\infty }dp\;\frac{p^{3}}{\gamma }%
\int_{-1}^{1}d\xi _{0}\;H\left( \left| \xi _{0}\right| -\xi _{0T}\right)
\;\xi _{0}f_{0}^{(0)}\left( p,\xi _{0},\psi \right) \\
&&\frac{1}{\overline{q}\left( \psi \right) }\int_{0}^{2\pi }\frac{d\theta }{%
2\pi }\frac{1}{\left| \widehat{\psi }\cdot \widehat{r}\right| }\frac{r}{R}%
\frac{B_{T}}{B_{P}}
\end{eqnarray}

We recognize the expression of the safety factor (\ref{qfactor}) so that 
\begin{equation}
\left\langle J_{\Vert }\right\rangle _{\phi }^{0}\left( \psi \right) =\frac{%
2\pi q_{e}}{m}\frac{q\left( \psi \right) }{\overline{q}\left( \psi \right) }%
\int_{0}^{\infty }dp\;\frac{p^{3}}{\gamma }\int_{-1}^{1}d\xi _{0}\;H\left(
\left| \xi _{0}\right| -\xi _{0T}\right) \;\xi _{0}f_{0}^{(0)}\left( p,\xi
_{0},\psi \right) \;
\end{equation}

\paragraph{Case of circular concentric flux-surfaces}

In that case, we showed in (\ref{qcirc}) that the safety factor is 
\begin{equation}
q\left( r\right) =\frac{\epsilon }{\sqrt{1-\epsilon ^{2}}}\frac{B_{T}}{B_{P}}
\label{qcirc}
\end{equation}
with $\epsilon =r/R_{0}$.

In addition, $\overline{q}\left( \psi \right) $ becomes 
\begin{eqnarray}
\overline{q}\left( \psi \right) &=&\int_{0}^{2\pi }\frac{d\theta }{2\pi }%
\frac{r}{R}\frac{B_{\min }\left( \psi \right) }{B_{P}} \\
&=&\int_{0}^{2\pi }\frac{d\theta }{2\pi }\frac{r}{R}\frac{B_{\min }\left(
\psi \right) }{B}\frac{B}{B_{P}} \\
&=&\epsilon \frac{B}{B_{P}}\frac{R_{0}}{R_{\max }} \\
&=&\frac{\epsilon }{1+\epsilon }\frac{B}{B_{P}}
\end{eqnarray}

We have then 
\begin{equation}
\frac{q\left( \psi \right) }{\overline{q}\left( \psi \right) }=\sqrt{\frac{
1+\epsilon }{1-\epsilon }}\frac{B_{T}}{B}
\end{equation}

In the case when $B_{T}\gg B_{P}$, we retrieve the bounce-averaged
coefficient $s^{*}$.

\subsubsection{Drift Kinetic Equation}

When we consider the first order distribution function, we have $f_{1}=%
\widetilde{f}+g$, where $g$ is constant along a field line, and therefore
its contribution has the same expression as for $f_{0}$. However, $%
\widetilde{f}$ has an explicit dependence upon $\theta $, which is given by (%
\ref{ftilde}) 
\begin{equation}
\widetilde{f}(p,\xi ,\psi ,\theta )=\frac{\xi (\psi ,\theta ,\xi _{0})}{\Psi
(\psi ,\theta )\xi _{0}}\widetilde{f}^{(0)}(p,\xi _{0},\psi )
\label{ftilde2}
\end{equation}
where 
\begin{equation}
\xi _{0}=\sigma \sqrt{1-\frac{1}{\Psi \left( \psi ,\theta \right) }\left(
1-\xi ^{2}\right) }
\end{equation}

Consequently, we find 
\begin{eqnarray}
\widetilde{J}_{\Vert }^{1}\left( \psi ,\theta \right) &=&2\pi
q_{e}\int_{0}^{\infty }p^{2}dp\int_{-1}^{1}d\xi \;\frac{p\xi }{\gamma m}%
\widetilde{f}\left( p,\xi ,\psi ,\theta \right)  \nonumber \\
&=&2\pi q_{e}\int_{0}^{\infty }p^{2}dp\int_{-1}^{1}d\xi \;\frac{p\xi }{%
\gamma m}\frac{\xi }{\Psi (\psi ,\theta )\xi _{0}}\widetilde{f}^{(0)}(p,\xi
_{0},\psi )  \nonumber \\
&=&2\pi q_{e}\int_{0}^{\infty }p^{2}dp\int_{-1}^{1}d\xi _{0}\frac{\xi }{\xi
_{0}}H\left( \left| \xi _{0}\right| -\sqrt{1-\frac{1}{\Psi \left( \psi
,\theta \right) }}\right) \;\frac{p\xi _{0}}{\gamma m}\widetilde{f}%
^{(0)}\left( p,\xi _{0},\psi \right)  \label{jpar}
\end{eqnarray}
where the condition 
\begin{equation}
\left| \xi _{0}\right| \geq \sqrt{1-\frac{1}{\Psi \left( \psi ,\theta
\right) }}  \label{xi0cond1}
\end{equation}
results from the equation (\ref{xi0}) and means that only the particle who
reach the position $\theta $ must be considered.

Therefore, the flux-surface averaged current density contribution from $%
\widetilde{f}$ 
\begin{equation}
\left\langle \widetilde{J}_{\Vert }\right\rangle _{\phi }^{1}\left( \psi
\right) =\frac{1}{\overline{q}\left( \psi \right) }\int_{0}^{2\pi }\frac{%
d\theta }{2\pi }\frac{1}{\left| \widehat{\psi }\cdot \widehat{r}\right| }%
\frac{r}{R}\frac{B_{T}}{B_{P}}\frac{\widetilde{J}_{\Vert }^{1}\left( \psi
,\theta \right) }{\Psi \left( \psi ,\theta \right) }
\end{equation}
becomes 
\begin{eqnarray}
\left\langle \widetilde{J}_{\Vert }\right\rangle _{\phi }^{1}\left( \psi
\right) &=&\frac{2\pi q_{e}}{m}\int_{0}^{\infty }dp\;\frac{p^{3}}{\gamma }%
\frac{1}{\overline{q}\left( \psi \right) }\int_{0}^{2\pi }\frac{d\theta }{%
2\pi }\frac{1}{\left| \widehat{\psi }\cdot \widehat{r}\right| }\frac{r}{R}%
\frac{B_{T}}{B_{P}}\frac{1}{\Psi \left( \psi ,\theta \right) } \\
&&\int_{-1}^{1}d\xi _{0}\frac{\xi }{\xi _{0}}H\left( \left| \xi _{0}\right| -%
\sqrt{1-\frac{1}{\Psi \left( \psi ,\theta \right) }}\right) \;\xi _{0}%
\widetilde{f}^{(0)}\left( p,\xi _{0},\psi \right)
\end{eqnarray}

Note that the condition (\ref{xi0cond1}) is equivalent to 
\begin{equation}
\theta _{\min }\left( \psi ,\xi _{0}\right) \leq \theta \leq \theta _{\max
}\left( \psi ,\xi _{0}\right)
\end{equation}
so that, permuting the integrals over $\theta $ and $\xi _{0}$, we find 
\begin{eqnarray}
\left\langle \widetilde{J}_{\Vert }\right\rangle _{\phi }^{1}\left( \psi
\right) &=&\frac{2\pi q_{e}}{m}\int_{0}^{\infty }dp\;\frac{p^{3}}{\gamma }%
\int_{-1}^{1}d\xi _{0}\;\xi _{0}\widetilde{f}^{(0)}\left( p,\xi _{0},\psi
\right)  \nonumber \\
&&\frac{1}{\overline{q}\left( \psi \right) }\int_{\theta _{\min }}^{\theta
_{\max }}\frac{d\theta }{2\pi }\frac{1}{\left| \widehat{\psi }\cdot \widehat{%
r}\right| }\frac{r}{R}\frac{B_{T}}{B_{P}}\frac{1}{\Psi \left( \psi ,\theta
\right) }\frac{\xi }{\xi _{0}}  \label{J1tilde}
\end{eqnarray}

We have 
\begin{equation}
\frac{r}{R}\frac{B_{T}}{B_{P}}\frac{1}{\Psi \left( \psi ,\theta \right) }=%
\frac{R_{0}I\left( \psi \right) }{R^{2}\left( \psi ,\theta _{0}\right)
B_{\min }\left( \psi \right) }\frac{r}{R_{0}}\frac{B}{B_{P}}\left[ \frac{%
R\left( \psi ,\theta _{0}\right) }{R\left( \psi ,\theta \right) }\right]
^{2}\Psi ^{-2}\left( \psi ,\theta \right)
\end{equation}

Then, noting the the integrand in (\ref{J1tilde}) is independent of $\sigma $%
, so that the sum over $\sigma $ for trapped particles can be added, we
obtain 
\begin{eqnarray}
\left\langle \widetilde{J}_{\Vert }\right\rangle _{\phi }^{1}\left( \psi
\right) &=&\frac{2\pi q_{e}}{m}\int_{0}^{\infty }dp\;\frac{p^{3}}{\gamma }%
\int_{-1}^{1}d\xi _{0}\;\xi _{0}\widetilde{f}^{(0)}\left( p,\xi _{0},\psi
\right) \frac{1}{\overline{q}\left( \psi \right) }\frac{R_{0}I\left( \psi
\right) }{R^{2}\left( \psi ,\theta _{0}\right) B_{\min }\left( \psi \right) }
\\
&&\left[ \frac{1}{2}\sum_{\sigma }\right] _{T}\int_{\theta _{\min }}^{\theta
_{\max }}\frac{d\theta }{2\pi }\frac{1}{\left| \widehat{\psi }\cdot \widehat{%
r}\right| }\frac{r}{R_{0}}\frac{B}{B_{P}}\frac{\xi _{0}}{\xi }\left[ \frac{%
R\left( \psi ,\theta _{0}\right) }{R\left( \psi ,\theta \right) }\right]
^{2}\Psi ^{-2}\left( \psi ,\theta \right) \left[ \frac{\xi }{\xi _{0}}%
\right] ^{2}
\end{eqnarray}

We recognize the expression of a bounce coefficient as defined in (\ref{lkm}%
) so that we get finally 
\begin{equation}
\left\langle \widetilde{J}_{\Vert }\right\rangle _{\phi }^{1}\left( \psi
\right) =\frac{2\pi q_{e}}{m}\frac{\widetilde{q}\left( \psi \right) }{%
\overline{q}\left( \psi \right) }\frac{R_{0}}{R\left( \psi ,\theta
_{0}\right) }\frac{B_{T}\left( \psi ,\theta _{0}\right) }{B_{\min }\left(
\psi \right) }\int_{0}^{\infty }dp\;\frac{p^{3}}{\gamma }\int_{-1}^{1}d\xi
_{0}\;\lambda _{2,-2,2}\xi _{0}\widetilde{f}^{(0)}\left( p,\xi _{0},\psi
\right)
\end{equation}
with 
\begin{equation}
\lambda _{2,-2,2}=\lambda \left\{ \left( \frac{\xi \left( \psi ,\theta ,\xi
_{0}\right) }{\xi _{0}}\right) ^{2}\Psi ^{-2}\left( \psi ,\theta \right)
\left( \frac{R\left( \psi ,\theta _{0}\right) }{R\left( \psi ,\theta \right) 
}\right) ^{2}\right\}
\end{equation}

\paragraph{Case of circular concentric flux-surfaces}

In that case, we showed in (\ref{qtilde2}) that $\widetilde{q}$ is 
\begin{equation}
\widetilde{q}\left( \psi \right) =\epsilon \frac{B}{B_{P}}  \label{qcirc}
\end{equation}
with $\epsilon =r/R_{0}$.

In addition, $\overline{q}\left( \psi \right) $ is 
\begin{equation}
\overline{q}\left( \psi \right) =\frac{\epsilon }{1+\epsilon }\frac{B}{B_{P}}
\end{equation}
and we have 
\begin{eqnarray}
\theta _{0} &=&0 \\
R\left( \psi ,\theta _{0}\right) &=&R_{\max }=R_{0}\left( 1+\epsilon \right)
\end{eqnarray}

We have then 
\begin{equation}
\frac{q\left( \psi \right) }{\overline{q}\left( \psi \right) }\frac{R_{0}}{%
R\left( \psi ,\theta _{0}\right) }\frac{B_{T}\left( \psi ,\theta _{0}\right) 
}{B_{\min }\left( \psi \right) }=\frac{B_{T}}{B}
\end{equation}

Also, in this case, 
\begin{equation}
\Psi \left( \psi ,\theta \right) =\frac{R\left( \psi ,\theta _{0}\right) }{%
R\left( \psi ,\theta \right) }
\end{equation}
so that 
\begin{equation}
\lambda _{2,-2,2}=\lambda _{2,0,0}=\lambda \left\{ \frac{\xi ^{2}}{\xi
_{0}^{2}}\right\} =\widetilde{s}^{*}
\end{equation}
where $\widetilde{s}^{*}$ has been mentionned in the litterature.

\subsection{Power Density Associated with a Flux}

\subsubsection{Definition}

The kinetic energy associated with a relativistic electron of momentum $p$
is 
\begin{equation}
\varepsilon =m_{e}c^{2}(\gamma -1)  \label{enelec}
\end{equation}

Then, the local energy density of electrons is 
\begin{equation}
{\cal E}\left( {\bf x}\right) =\int d^{3}p\ m_{e}c^{2}(\gamma -1)f({\bf x},%
{\bf p})
\end{equation}

The density of power absorbed through the process ${\cal O}$ is 
\begin{equation}
P_{abs}^{{\cal O}}\left( {\bf x}\right) =\left. \frac{\partial {\cal E}}{%
\partial t}\right| _{{\cal O}}=\int d^{3}p\ m_{e}c^{2}\left( \gamma
-1\right) \left. \frac{\partial f({\bf x},{\bf p})}{\partial t}\right| _{%
{\cal O}}
\end{equation}

When the operator is described in conservative form, as the divergence of a
flux

\begin{equation}
\left. \frac{\partial f}{\partial t}\right| _{{\cal O}}=-\nabla _{{\bf p}
}\cdot {\bf S}_{{\bf p}}^{{\cal O}}=-\frac{1}{p^{2}}\frac{\partial }{%
\partial p}\left( p^{2}S_{p}^{{\cal O}}\right) +\frac{1}{p}\frac{\partial }{%
\partial \xi }\left( \sqrt{1-\xi ^{2}}S_{\xi }^{{\cal O}}\right)
\end{equation}
then the power density becomes

\begin{equation}
P_{abs}^{{\cal O}}=-2\pi m_{e}c^{2}\int_{0}^{\infty }p^{2}dp\;\left( \gamma
-1\right) \int_{-1}^{+1}d\xi \left[ \frac{1}{p^{2}}\frac{\partial }{\partial
p}\left( p^{2}S_{p}^{{\cal O}}\right) -\frac{1}{p}\frac{\partial }{\partial
\xi }\left( \sqrt{1-\xi ^{2}}S_{\xi }^{{\cal O}}\right) \right]  \label{pabs}
\end{equation}

The integration of the $S_{\xi }^{{\cal O}}$ term gives no contribution,
since the particle energy is function of $p$ only: 
\begin{equation}
\int_{-1}^{+1}d\xi \;\frac{\partial }{\partial \xi }\left( \sqrt{1-\xi ^{2}}
S_{\xi }^{{\cal O}}\right) =\left[ \sqrt{1-\xi ^{2}}S_{\xi }^{{\cal O}%
}\right] _{-1}^{+1}=0
\end{equation}
and the equation (\ref{pabs}) reduces to 
\begin{equation}
P_{abs}^{{\cal O}}=-2\pi m_{e}c^{2}\int_{-1}^{+1}d\xi \int_{0}^{\infty
}\left( \gamma -1\right) \frac{\partial }{\partial p}\left( p^{2}S_{p}^{%
{\cal O}}\right) dp
\end{equation}

Integrating by parts, we get 
\begin{equation}
P_{abs}^{{\cal O}}=-2\pi m_{e}c^{2}\int_{-1}^{+1}d\xi \left( \left[ \left(
\gamma -1\right) p^{2}S_{p}^{{\cal O}}\right] _{0}^{\infty
}-\int_{0}^{\infty }\frac{d\gamma }{dp}p^{2}S_{p}^{{\cal O}}dp\right)
\label{pabsb}
\end{equation}

Assuming that $\lim_{p\rightarrow \infty }p^{2}S_{p}^{{\cal O}}=0$, and
using 
\begin{equation}
\frac{d\gamma }{dp}=\frac{p}{\gamma m_{e}^{2}c^{2}}
\end{equation}
the equation (\ref{pabsb}) reduces to 
\begin{equation}
P_{abs}^{{\cal O}}\left( \psi ,\theta \right) =2\pi \int_{-1}^{+1}d\xi
\int_{0}^{\infty }dp\;\frac{p^{3}}{\gamma m_{e}}S_{p}^{{\cal O}}
\end{equation}

\subsubsection{Flux-Surface Averaging}

The flux-surface averaged power density is given by () 
\begin{equation}
\left\langle P_{abs}^{{\cal O}}\right\rangle _{V}\left( \psi \right) =\frac{1%
}{\widehat{q}\left( \psi \right) }\int_{0}^{2\pi }\frac{d\theta }{2\pi }%
\frac{1}{\left| \widehat{\psi }\cdot \widehat{r}\right| }\frac{r}{R_{0}}%
\frac{B_{\min }\left( \psi \right) }{B_{P}}P_{abs}^{{\cal O}}\left( \psi
,\theta \right)
\end{equation}
which becomes 
\begin{equation}
\left\langle P_{abs}^{{\cal O}}\right\rangle _{V}\left( \psi \right) =2\pi
\int_{0}^{\infty }dp\;\frac{p^{3}}{\gamma m_{e}}\frac{1}{\widehat{q}\left(
\psi \right) }\int_{0}^{2\pi }\frac{d\theta }{2\pi }\frac{1}{\left| \widehat{%
\psi }\cdot \widehat{r}\right| }\frac{r}{R_{0}}\frac{B_{\min }\left( \psi
\right) }{B_{P}}\int_{-1}^{+1}d\xi \;S_{p}^{{\cal O}}
\end{equation}

The sum over $\sigma $ for trapped electrons can be added, using 
\begin{eqnarray*}
\int_{-1}^{1}d\xi \left[ \frac{1}{2}\sum_{\sigma =\pm 1}\right] _{T}S_{p}^{%
{\cal O}} &=&\int_{-1}^{-\xi _{T}}d\xi S_{p}^{{\cal O}}+\int_{\xi
_{T}}^{1}d\xi S_{p}^{{\cal O}}+\frac{1}{2}\int_{-\xi _{T}}^{\xi _{T}}d\xi
\sum_{\sigma =\pm 1}S_{p}^{{\cal O}} \\
&=&\int_{-1}^{-\xi _{T}}d\xi S_{p}^{{\cal O}}+\int_{\xi _{T}}^{1}d\xi S_{p}^{%
{\cal O}}+\frac{1}{2}\int_{-\xi _{T}}^{\xi _{T}}d\xi \left[ S_{p}^{{\cal O}%
}(\xi )+S_{p}^{{\cal O}}(-\xi )\right] \\
&=&\int_{-1}^{-\xi _{T}}d\xi S_{p}^{{\cal O}}+\int_{\xi _{T}}^{1}d\xi S_{p}^{%
{\cal O}}+\int_{-\xi _{T}}^{\xi _{T}}d\xi S_{p}^{{\cal O}}(\xi ) \\
&=&\int_{-1}^{1}d\xi S_{p}^{{\cal O}}
\end{eqnarray*}
where the trapping condition evaluated at the location $\theta $ is 
\begin{equation}
\left| \xi \right| <\xi _{T}=\sqrt{1-\frac{B\left( \psi ,\theta \right) }{%
B_{\max }\left( \psi \right) }}
\end{equation}

Using $\xi d\xi =\Psi \xi _{0}d\xi _{0}$ with the condition (\ref{xi0cond})
on $\xi _{0}$%
\begin{equation}
\left| \xi _{0}\right| \geq \sqrt{1-\frac{1}{\Psi \left( \psi ,\theta
\right) }}  \label{xi0cond2}
\end{equation}
we get that 
\begin{equation}
\int_{-1}^{+1}d\xi =\int_{-1}^{+1}d\xi _{0}\;\frac{\Psi \left( \psi ,\theta
\right) \xi _{0}}{\xi }H\left( \left| \xi _{0}\right| -\sqrt{1-\frac{1}{\Psi
\left( \psi ,\theta \right) }}\right)
\end{equation}

Note that the condition (\ref{xi0cond2}) is equivalent to 
\begin{equation}
\theta _{\min }\left( \psi ,\xi _{0}\right) \leq \theta \leq \theta _{\max
}\left( \psi ,\xi _{0}\right)  \label{thetacond}
\end{equation}
so that, permuting the integrals over $\theta $ and $\xi _{0}$, we find 
\begin{eqnarray}
\left\langle P_{abs}^{{\cal O}}\right\rangle _{V}\left( \psi \right) &=&2\pi
\int_{0}^{\infty }dp\;\frac{p^{3}}{\gamma m_{e}}\int_{-1}^{+1}d\xi _{0} \\
&&\frac{1}{\widehat{q}\left( \psi \right) }\left[ \frac{1}{2}\sum_{\sigma
=\pm 1}\right] _{T}\int_{\theta _{\min }}^{\theta _{\max }}\frac{d\theta }{%
2\pi }\frac{1}{\left| \widehat{\psi }\cdot \widehat{r}\right| }\frac{r}{R_{0}%
}\frac{B}{B_{P}}\frac{\xi _{0}}{\xi }S_{p}^{{\cal O}}
\end{eqnarray}

We see that the bounce-averaging of the fluxes appears naturally, so that we
can rewrite 
\begin{equation}
\left\langle P_{abs}^{{\cal O}}\right\rangle _{V}\left( \psi \right) =2\pi 
\frac{\widetilde{q}\left( \psi \right) }{\widehat{q}\left( \psi \right) }%
\int_{0}^{\infty }dp\;\frac{p^{3}}{\gamma m_{e}}\int_{-1}^{+1}d\xi
_{0}\lambda \left\{ S_{p}^{{\cal O}}\right\}
\end{equation}

\paragraph{Case of circular concentric flux-surfaces}

In that case, we showed in (\ref{qcirc}) that the coefficient $\widetilde{q}$
is 
\begin{equation}
\widetilde{q}\left( \psi \right) =\epsilon \frac{B}{B_{P}}  \label{qtilde2}
\end{equation}
with $\epsilon =r/R_{0}$.

In addition, $\widehat{q}\left( \psi \right) $ becomes 
\begin{eqnarray}
\widehat{q}\left( \psi \right) &=&\int_{0}^{2\pi }\frac{d\theta }{2\pi }%
\frac{r}{R_{0}}\frac{B_{\min }\left( \psi \right) }{B_{P}} \\
&=&\epsilon \frac{B}{B_{P}}\frac{R_{0}}{R_{\max }}\int_{0}^{2\pi }\frac{%
d\theta }{2\pi }\frac{R}{R_{0}} \\
&=&\frac{\epsilon }{1+\epsilon }\frac{B}{B_{P}}
\end{eqnarray}

We have then 
\begin{equation}
\frac{\widetilde{q}\left( \psi \right) }{\widehat{q}\left( \psi \right) }%
=1+\epsilon 
\end{equation}

\subsection{Stream Function for Momentum Space fluxes}

When transport in configuration space is ignored, and a steady-state is
assumed, the Fokker-Planck equation (\ref{consform}) reduces to the
conservative equation (\ref{momentumfpe})
\begin{equation}
\nabla _{{\bf p}}\cdot {\bf S}_{{\bf p}}=0
\end{equation}

Because ${\bf S}_{{\bf p}}$ is a divergence-free field vector, it can be
expressed as the curl of a stream function
\begin{equation}
{\bf S}_{{\bf p}}=\nabla \times {\bf T}_{{\bf p}}
\end{equation}

The expression of a curl in momentum space $\left( p,\xi ,\varphi \right) $
gives (\ref{pxicurl})
\begin{eqnarray}
S_{p} &=&\frac{1}{p}\frac{\partial }{\partial \xi }\left( \sqrt{1-\xi ^{2}}%
T_{\varphi }\right) +\frac{1}{p\sqrt{1-\xi ^{2}}}\frac{\partial T_{\xi }}{%
\partial \varphi } \\
S_{\xi } &=&\frac{1}{p}\frac{\partial }{\partial p}\left( pT_{\varphi
}\right) -\frac{1}{p\sqrt{1-\xi ^{2}}}\frac{\partial T_{p}}{\partial \varphi 
} \\
S_{\varphi } &=&-\frac{1}{p}\frac{\partial }{\partial p}\left( pT_{\xi
}\right) -\frac{\sqrt{1-\xi ^{2}}}{p}\frac{\partial T_{p}}{\partial \xi }
\end{eqnarray}

Because $S_{\varphi }=0$, we can choose $T_{\xi }=T_{p}=0$, which leads to
\begin{eqnarray}
S_{p} &=&\frac{1}{p}\frac{\partial }{\partial \xi }\left( \sqrt{1-\xi ^{2}}%
T_{\varphi }\right)  \\
S_{\xi } &=&\frac{1}{p}\frac{\partial }{\partial p}\left( pT_{\varphi
}\right) 
\end{eqnarray}
and we can rewrite
\begin{equation}
{\bf S}_{{\bf p}}=\nabla \times T_{\varphi }\widehat{\varphi }
\end{equation}

In order to give a physical meaning to $T_{\varphi }\left( p,\xi ,\psi
\right) $, we define formally 
\begin{equation}
T_{\varphi }\left( p,\xi ,\psi \right) =K\left( p,\xi ,\psi \right) A\left(
p,\xi ,\psi \right) 
\end{equation}
where the function $A\left( p,\xi \right) $ is such that the flux of
electrons between two contours $A_{1}$ and $A_{2}$ is equal to $n\left( \psi
\right) $ $\left( A_{2}-A_{1}\right) $. Lets consider a path $\gamma _{12}$
between the contours $A_{1}$ and $A_{2}$. The total flux of electrons
through this path, which is in fact a surface, given the rotational symmetry
in $\varphi $, is given by
\begin{eqnarray}
\Gamma _{12} &=&\iint_{S_{12}}dS\;{\bf S}_{{\bf p}}\cdot \widehat{n} \\
&=&\iint_{S_{12}}d{\bf S}\cdot \nabla \times T_{\varphi }\widehat{\varphi }
\\
&=&\oint_{{\cal C}_{12}}T_{\varphi }\;d{\bf l}\cdot \widehat{\varphi }
\end{eqnarray}

By rotational symmetry in $\varphi $, and using (\ref{dlphipxi}), we get that
\begin{equation}
\Gamma _{12}=2\pi p_{2}\sqrt{1-\xi _{2}^{2}}T_{\varphi 2}\;-2\pi p_{1}\sqrt{%
1-\xi _{1}^{2}}T_{\varphi 1}
\end{equation}
If we define
\begin{equation}
K\left( p,\xi ,\psi \right) \equiv \frac{n\left( \psi \right) }{2\pi p\sqrt{%
1-\xi ^{2}}}
\end{equation}
we get
\begin{equation}
\Gamma _{12}=n\left( \psi \right) \left( A_{2}\;-A_{1}\right) 
\end{equation}
and therefore we obtain that the total flux between the contours $A_{1}$ and 
$A_{2}$ is equal to $n\left( \psi \right) $ $\left( A_{2}-A_{1}\right) $. We
call $A\left( p,\xi ,\psi \right) $ the stream function, and we get finally
\begin{eqnarray}
S_{p} &=&\frac{n\left( \psi \right) }{2\pi p^{2}}\frac{\partial A}{\partial %
\xi } \\
S_{\xi } &=&\frac{n\left( \psi \right) }{2\pi p\sqrt{1-\xi ^{2}}}\frac{%
\partial A}{\partial p}
\end{eqnarray}

Since there are no fluxes accross the internal boundaries in the momentum
space, this boundary coincide with a contour $A$, and therefore we can
arbitrarily set thsi value to $0$:
\begin{equation}
A\left( 0,\xi \right) =A\left( p,\pm 1\right) =0
\end{equation}
Then $A$ can be calculated by any of the integrals
\begin{eqnarray}
A\left( p,\xi ,\psi \right)  &=&\frac{2\pi p^{2}}{n\left( \psi \right) }%
\int_{-1}^{\xi }d\xi \;S_{p}=\frac{2\pi p^{2}}{n\left( \psi \right) }%
\int_{1}^{\xi }d\xi ^{\prime }\;S_{p} \\
A\left( p,\xi ,\psi \right)  &=&\frac{2\pi \sqrt{1-\xi ^{2}}}{n\left( \psi
\right) }\int_{0}^{p}p^{\prime }dp^{\prime }\;S_{\xi }
\end{eqnarray}

However, $A\left( p,\xi ,\psi \right) $ remains a function of $\xi $, which
depends upon $\theta $. Starting from the bounce-averaged fluxes, it would
be interesting to compute a function $A^{\left( 0\right) }\left( p,\xi
_{0},\psi \right) $, such that 
\begin{eqnarray}
A^{\left( 0\right) }\left( 0,\xi _{0}\right)  &=&\overline{A}\left( p,\pm %
1\right) =0  \nonumber \\
S_{p}^{\left( 0\right) } &=&\frac{n\left( \psi \right) }{2\pi p^{2}}\frac{%
\partial A^{\left( 0\right) }}{\partial \xi _{0}}  \label{A0} \\
S_{\xi }^{\left( 0\right) } &=&\frac{n\left( \psi \right) }{2\pi p\sqrt{%
1-\xi _{0}^{2}}}\frac{\partial A^{\left( 0\right) }}{\partial p}  \nonumber
\end{eqnarray}

We first need to demonstrate the existence of such a function. Starting from 
$S_{p}^{\left( 0\right) }$,
\begin{eqnarray}
A^{\left( 0\right) }\left( p,\xi _{0},\psi \right)  &=&\frac{2\pi p^{2}}{
n\left( \psi \right) }\int_{-1}^{\xi _{0}}d\xi _{0}^{\prime }\;\left\{
S_{p}\right\}  \\
&=&\frac{2\pi p^{2}}{n\left( \psi \right) }\int_{-1}^{\xi _{0}}d\xi
_{0}^{\prime }\;\frac{1}{\lambda \widetilde{q}}\left[ \frac{1}{2}%
\sum_{\sigma }\right] _{T}\int_{\theta _{\min }}^{\theta _{\max }}\frac{%
d\theta }{2\pi }\frac{1}{\left| \widehat{\psi }\cdot \widehat{r}\right| }%
\frac{r}{R_{0}}\frac{B}{B_{P}}\frac{\xi _{0}^{\prime }}{\xi ^{\prime }}S_{p}
\\
&=&\int_{-1}^{\xi _{0}}d\xi _{0}^{\prime }\;\frac{1}{\lambda \widetilde{q}}%
\left[ \frac{1}{2}\sum_{\sigma }\right] _{T}\int_{0}^{2\pi }\frac{d\theta }{2%
\pi }H\left( B_{b}-B\right) \frac{1}{\left| \widehat{\psi }\cdot \widehat{r}%
\right| }\frac{r}{R_{0}}\frac{B}{B_{P}}\frac{\xi _{0}^{\prime }}{\xi
^{\prime }}\frac{\partial A}{\partial \xi ^{\prime }} \\
&=&\;\frac{\sigma }{\lambda \widetilde{q}}\left[ \frac{1}{2}\sum_{\sigma
}\right] _{T}\int_{0}^{2\pi }\frac{d\theta }{2\pi }\frac{1}{\left| \widehat{
\psi }\cdot \widehat{r}\right| }\frac{r}{R_{0}}\frac{B}{B_{P}}\int_{-1}^{\xi
_{0}}d\xi _{0}^{\prime }H\left( \left| \xi _{0}\right| -\sqrt{1-\frac{1}{%
\Psi }}\right) \sigma \frac{\xi _{0}^{\prime }}{\xi ^{\prime }}\frac{%
\partial A}{\partial \xi ^{\prime }} \\
&=&\;\frac{\sigma }{\lambda \widetilde{q}}\left[ \frac{1}{2}\sum_{\sigma
}\right] _{T}\int_{0}^{2\pi }\frac{d\theta }{2\pi }\frac{1}{\left| \widehat{%
\psi }\cdot \widehat{r}\right| }\frac{r}{R_{0}}\frac{B}{B_{P}}\frac{\sigma }{%
\Psi }\int_{-1}^{\xi }d\xi ^{\prime }\frac{\partial A}{\partial \xi ^{\prime
}} \\
&=&\;\frac{\sigma }{\lambda \widetilde{q}}\left[ \frac{1}{2}\sum_{\sigma
}\right] _{T}\int_{0}^{2\pi }\frac{d\theta }{2\pi }\frac{1}{\left| \widehat{%
\psi }\cdot \widehat{r}\right| }\frac{r}{R_{0}}\frac{B_{\min }\left( \psi
\right) }{B_{P}}\sigma A \\
&=&\sigma \frac{\widehat{q}\left( \psi \right) }{\lambda \widetilde{q}}%
\left[ \frac{1}{2}\sum_{\sigma }\right] _{T}\left\langle \sigma
A\right\rangle _{V}
\end{eqnarray}
where we used that
\begin{equation}
\theta _{\min }\leq \theta \leq \theta _{\max }\Leftrightarrow B\leq %
B_{b}\Leftrightarrow \sqrt{1-\frac{1}{\Psi }}\leq \left| \xi _{0}\right| 
\end{equation}

Now, starting from $S_{\xi }^{\left( 0\right) }$, we have
\begin{eqnarray}
A^{\left( 0\right) }\left( p,\xi _{0},\psi \right)  &=&\frac{2\pi \sqrt{%
1-\xi _{0}^{2}}}{n\left( \psi \right) }\int_{0}^{p}p^{\prime }dp^{\prime
}\;\sigma \left\{ \frac{\sigma \xi }{\sqrt{\Psi }\xi _{0}}S_{\xi }\right\} 
\\
&=&\frac{2\pi \sqrt{1-\xi _{0}^{2}}}{n\left( \psi \right) }%
\int_{0}^{p}p^{\prime }dp^{\prime }\;\sigma \frac{1}{\lambda \widetilde{q}}%
\left[ \frac{1}{2}\sum_{\sigma }\right] _{T}\int_{\theta _{\min }}^{\theta
_{\max }}\frac{d\theta }{2\pi }\frac{1}{\left| \widehat{\psi }\cdot \widehat{%
r}\right| }\frac{r}{R_{0}}\frac{B}{B_{P}}\frac{\sigma }{\sqrt{\Psi }}S_{\xi }
\\
&=&\int_{0}^{p}dp^{\prime }\;\sigma \frac{1}{\lambda \widetilde{q}}\left[ 
\frac{1}{2}\sum_{\sigma }\right] _{T}\int_{\theta _{\min }}^{\theta _{\max }}%
\frac{d\theta }{2\pi }\frac{1}{\left| \widehat{\psi }\cdot \widehat{r}%
\right| }\frac{r}{R_{0}}\frac{B}{B_{P}}\frac{\sigma }{\sqrt{\Psi }}\frac{%
\sqrt{1-\xi _{0}^{2}}}{\sqrt{1-\xi ^{2}}}\frac{\partial A}{\partial %
p^{\prime }} \\
&=&\;\frac{\sigma }{\lambda \widetilde{q}}\left[ \frac{1}{2}\sum_{\sigma
}\right] _{T}\int_{\theta _{\min }}^{\theta _{\max }}\frac{d\theta }{2\pi }%
\frac{1}{\left| \widehat{\psi }\cdot \widehat{r}\right| }\frac{r}{R_{0}}%
\frac{B}{B_{P}}\frac{\sigma }{\Psi }\int_{0}^{p}dp^{\prime }\frac{\partial A%
}{\partial p^{\prime }} \\
&=&\;\frac{\sigma }{\lambda \widetilde{q}}\left[ \frac{1}{2}\sum_{\sigma
}\right] _{T}\int_{\theta _{\min }}^{\theta _{\max }}\frac{d\theta }{2\pi }%
\frac{1}{\left| \widehat{\psi }\cdot \widehat{r}\right| }\frac{r}{R_{0}}%
\frac{B_{\min }\left( \psi \right) }{B_{P}}\sigma A \\
&=&\sigma \frac{\widehat{q}\left( \psi \right) }{\lambda \widetilde{q}}%
\left[ \frac{1}{2}\sum_{\sigma }\right] _{T}\left\langle \sigma
A\right\rangle _{V}
\end{eqnarray}
and we find the same function $A^{\left( 0\right) }$. The existence of a
function $A^{\left( 0\right) }$ verifying (\ref{A0}) is therefore
demonstrated. We need now t demonstrate that $A^{\left( 0\right) }$
verifying (\ref{A0}) leads to the bounce-averaged FP equation (\ref{kinbav2}%
):
\begin{eqnarray}
\left\{ \left. \frac{\partial f}{\partial t}\right| _{{\bf p}}\right\}  &=&%
\frac{1}{p^{2}}\frac{\partial }{\partial p}\left( p^{2}S_{p}^{\left(
0\right) }\right) -\frac{1}{\lambda p}\frac{\partial }{\partial \xi _{0}}%
\left( \sqrt{1-\xi _{0}^{2}}\lambda S_{\xi _{0}}^{\left( 0\right) }\right) 
\label{kinbav2} \\
&=&\frac{1}{p^{2}}\frac{\partial }{\partial p}\left( p^{2}\frac{n\left( \psi
\right) }{2\pi p^{2}}\frac{\partial A^{\left( 0\right) }}{\partial \xi _{0}}%
\right) -\frac{1}{\lambda p}\frac{\partial }{\partial \xi _{0}}\left( \sqrt{%
1-\xi _{0}^{2}}\lambda \frac{n\left( \psi \right) }{2\pi p\sqrt{1-\xi
_{0}^{2}}}\frac{\partial A^{\left( 0\right) }}{\partial p}\right)  \\
&=&\frac{1}{\lambda p^{2}}\frac{\partial ^{2}}{\partial p\xi _{0}}\left[ 
\frac{\lambda n\left( \psi \right) A^{\left( 0\right) }}{2\pi }\right] -%
\frac{1}{\lambda p^{2}}\frac{\partial ^{2}}{\partial \xi _{0}\partial p}%
\left[ \frac{\lambda n\left( \psi \right) A^{\left( 0\right) }}{2\pi }%
\right]  \\
&=&0
\end{eqnarray}

In conclusion, we found a stream function verifying
\begin{equation}
A^{\left( 0\right) }\left( 0,\xi _{0}\right) =\overline{A}\left( p,\pm %
1\right) =0
\end{equation}
which leads to the bounce-averaged FP equation and which can be calculated
from the bounce-averaged fluxes by either
\begin{equation}
A^{\left( 0\right) }=\frac{2\pi p^{2}}{n\left( \psi \right) }\int_{-1}^{\xi
_{0}}d\xi _{0}^{\prime }\;S_{p}^{\left( 0\right) }=\frac{2\pi p^{2}}{n\left(
\psi \right) }\int_{1}^{\xi _{0}}d\xi _{0}^{\prime }\;S_{p}^{\left( 0\right)
}  \nonumber
\end{equation}
or
\begin{equation}
A^{\left( 0\right) }=\frac{2\pi \sqrt{1-\xi _{0}^{2}}}{n\left( \psi \right) }%
\int_{0}^{p}p^{\prime }dp^{\prime }\;S_{\xi }^{\left( 0\right) }
\end{equation}

\subsection{Fraction of trapped electrons}

The effective fraction of trapped electrons $f_{t}$ is one of the most
important quantities of the neoclassical transport theory, since the
parallel viscosity responsible for reduction of the Ohmic conductivity and
the bootstrap current are directly proportional to $f_{t}\left(
1-f_{t}\right) $ in the Hirschan-Sigmar formulation of the transport theory.
It is a purely kinetic quantity, that can only be estimated localy at the
common plasma location for all electrons trajectories i.e. at $B_{min}$. In
the thin banana width approximation, it is therefore a function of the flux
surface $\psi $%
\begin{equation}
f_{t}^{kin}\left( \psi \right) =\frac{2\pi \int_{0}^{\infty %
}p^{2}dp\int_{-\xi _{0T}}^{+\xi _{0T}}f^{(0)}\left( \psi ,p,\xi _{0}\right)
d\xi _{0}}{n_{e}\left( \psi \right) }  \label{kintrapfracgen}
\end{equation}
where $\xi _{0T}\left( \psi \right) $ is given by the relation (\ref{xi0t2})
and $n_{e}\left( \psi \right) $ is the flux-surface electron averaged
density. Recalling the expansion up to the first order
\begin{equation}
f^{(0)}\simeq f_{0}^{(0)}+f_{1}^{(0)}
\end{equation}
where $f_{1}^{(0)}=\widetilde{f}^{\left( 0\right) }+g^{\left( 0\right) },$
it turns out that
\begin{equation}
f_{t}^{kin}\left( \psi \right) \simeq f_{0,t}^{kin}\left( \psi \right) +%
\widetilde{f}_{t}^{kin}\left( \psi \right) +g_{t}^{kin}\left( \psi \right) 
\end{equation}
where
\begin{eqnarray}
f_{0,t}^{kin}\left( \psi \right)  &=&\frac{2\pi }{n_{e}\left( \psi \right) }%
\int_{0}^{\infty }p^{2}dp\int_{-\xi _{0T}}^{+\xi _{0T}}f_{0}^{(0)}\left(
\psi ,p,\xi _{0}\right) d\xi _{0}  \nonumber \\
\widetilde{f}_{t}^{kin}\left( \psi \right)  &=&\frac{2\pi }{n_{e}\left( \psi
\right) }\int_{0}^{\infty }p^{2}dp\int_{-\xi _{0T}}^{+\xi _{0T}}\widetilde{f}%
^{\left( 0\right) }\left( \psi ,p,\xi _{0}\right) d\xi _{0}  \nonumber \\
g_{t}^{kin}\left( \psi \right)  &=&\frac{2\pi }{n_{e}\left( \psi \right) }%
\int_{0}^{\infty }p^{2}dp\int_{-\xi _{0T}}^{+\xi _{0T}}g^{\left( 0\right)
}\left( \psi ,p,\xi _{0}\right) d\xi _{0}
\end{eqnarray}
gradient terms taking into account of the neoclassical corrections. Since $%
\widetilde{f}^{\left( 0\right) }$ is always antisymmetric in the trapped
region, $\widetilde{f}_{t}^{kin}\left( \psi \right) =0$ and
\begin{equation}
f_{t}^{kin}\left( \psi \right) \simeq f_{0,t}^{kin}\left( \psi \right)
+g_{t}^{kin}\left( \psi \right)   \label{trapfrackin}
\end{equation}

For a Maxwellian distribution function, since in addition $g_{M}^{\left(
0\right) }=0$ in the trapped region,
\begin{equation}
g_{t,M}^{kin}\left( \psi \right) =0
\end{equation}
and since $f_{0,M}^{(0)}$ is only a function of $p$ and $\psi ,$

\begin{equation}
f_{0,t,M}^{kin}\left( \psi \right) =\frac{\int_{-\xi _{0T}}^{+\xi _{0T}}d\xi
_{0}}{\int_{-1}^{+1}d\xi _{0}}=\xi _{0T}=\sqrt{1-\frac{B_{\min }\left( \psi
\right) }{B_{\max }\left( \psi \right) }}
\end{equation}
This expression is different from the one deduced from neoclassical theory,
whose general formulation is 
\begin{equation}
f_{t,M}^{neo}\left( \psi \right) =1-\frac{3}{4}\left\langle
h^{2}\right\rangle \int_{0}^{1}\frac{xdx}{\left\langle \sqrt{1-xh}%
\right\rangle }  \label{neotrapfracgenthermal}
\end{equation}
where $h=B/B_{max}$ and $B_{max}$ is the maximum value of the magnetic field 
$B$ along the particle trajectory. This difference likely arises from second
order terms that has been incorporated in the neoclassical transport theory
for thermal regimes, that are here neglected. However, when RF waves
strongly distort the electron distribution function, the first term $%
g_{t}^{kin}\left( \psi \right) $ may become more important and only the
expression (\ref{kintrapfracgen}) must be considered.

Regarding the numerical application, expression (\ref{kintrapfracgen}) is
more general for arbitrary shapes of the electron distribution function, and
involves quantities that are naturally determined by the Fokker-Planck
equation solver. However, it is interesting to compare with the general
formula (\ref{neotrapfracgenthermal}) in the thermal regime, in order to
validate both approaches.

Since $f_{t,M}^{neo}$ involves a cumbersome double integral that is
numerically demanding in the vicinity of the core region of the plasma,
several attempts have been done to obtain a simplified, though accurate,
expression that may be used in all types of equilibria, especially in high $%
\beta _{p}$ regime. So far, the best approach in the one performed by Y.R.
Lin-Liu and R. L. Miller (Phys. Plasma, 2 (1995) 1666) that consider upper
and lower limit boundaries of (\ref{neotrapfracgenthermal}) using Schwartz
inequality. Hence, 

\begin{equation}
f_{tl,M}^{neo}\left( \psi \right) \leq f_{t,M}^{neo}\left( \psi \right) \leq %
f_{tu,M}^{neo}\left( \psi \right) 
\end{equation}
with 
\begin{equation}
f_{tu,M}^{neo}\left( \psi \right) =1-\frac{\left\langle h^{2}\right\rangle }{%
\left\langle h\right\rangle ^{2}}\left[ 1-\sqrt{1-\left\langle
h\right\rangle }\left( 1+\frac{\left\langle h\right\rangle }{2}\right)
\right] 
\end{equation}
and
\begin{equation}
f_{tl,M}^{neo}\left( \psi \right) =1-\left\langle h^{2}\right\rangle
\left\langle \frac{1-\sqrt{1-h}\left( 1+\frac{h}{2}\right) }{h^{2}}%
\right\rangle 
\end{equation}

By expanding $f_{tu,M}^{neo}$ and $f_{tl,M}^{neo}$ in terms of moments of $h$%
, one obtains inifinite series of weak convergence for a simple practical
purpose, but the set of coefficients suggest that a linear combination of $%
f_{tu,M}^{neo}$ and $f_{tl,M}^{neo},$ may provide an accurate enough
estimate of $f_{t,M}^{neo}$ for nearly all plasma regimes. Therefore, 
\begin{equation}
\overline{f}_{t,M}^{neo}\approx \varpi f_{tu,M}^{neo}+(1-\varpi
)f_{tl,M}^{neo}  \label{intertrapfrac}
\end{equation}
where the usual value for $\varpi $ deduced from numerical simulations is $%
0.75$. 

According to the definition of the volume flux surface averaged, 
\begin{equation}
\left\langle h\right\rangle =\frac{1}{\widehat{q}\left( \psi \right) }\frac{%
B_{\min }\left( \psi \right) }{B_{\max }\left( \psi \right) }\int_{0}^{2\pi }%
\frac{d\theta }{2\pi }\frac{1}{\left| \widehat{\psi }\cdot \widehat{r}%
\right| }\frac{r}{R_{0}}\frac{B}{B_{P}}
\end{equation}
where
\begin{equation}
\widehat{q}\left( \psi \right) =\int_{0}^{2\pi }\frac{d\theta }{2\pi }\frac{1%
}{\left| \widehat{\psi }\cdot \widehat{r}\right| }\frac{r}{R_{0}}\frac{%
B_{\min }\left( \psi \right) }{B_{P}}
\end{equation}

\subsubsection{Case of circular concentric flux surfaces}

In that case, $\epsilon =r/R_{0},$ $\widehat{\psi }\cdot \widehat{r}=1,$ and
the ratio $B/B_{P}$ is a function of $r$ only. Therefore, since 
\begin{equation}
\frac{B_{\min }\left( r\right) }{B_{\max }\left( r\right) }=\frac{1-\epsilon 
}{1+\epsilon }
\end{equation}
It is then straightforward to evaluate $f_{t,M}^{kin}\left( r\right) $ which
is
\begin{equation}
f_{t,M}^{kin}\left( r\right) =\sqrt{\frac{2\epsilon }{1+\epsilon }}
\end{equation}
and in the limit of large aspect ratio, 
\begin{equation}
\lim_{\epsilon \rightarrow 0}f_{t,M}^{kin}\left( r\right) =\sqrt{2\epsilon }
\end{equation}
which is very close to the exact value $f_{t,M}^{neo}$ of the neoclassical
theory $\gamma \sqrt{\epsilon }$ with $\gamma \simeq 1.46\footnote{%
The exact analytical result according to FL. Hinton and R.D. Hazeltine (Rev.
Mod. Phys. \textbf{48}, 239 (1976)), gives for $\gamma ,$%
\[
\gamma \equiv \frac{3\sqrt{2}}{2}\left[ 1+\int 01\frac{dx}{x^{2}}\left( 1-%
\frac{\pi }{2E\left( x\right) }\right) \right] \cong 1.46
\]
}.$

For the determination of $f_{tl}^{neo}$ and $f_{tu}^{neo},$ 
\begin{equation}
\left\langle h\right\rangle =\frac{1}{\widehat{q}\left( r\right) }\frac{%
1-\epsilon }{1+\epsilon }\epsilon \frac{B}{B_{P}}
\end{equation}
and 
\begin{equation}
\widehat{q}\left( r\right) =\frac{\epsilon }{1+\epsilon }\frac{B}{B_{P}}
\end{equation}
\begin{equation}
\left\langle h\right\rangle =1-\epsilon 
\end{equation}

In the same way, 
\begin{eqnarray}
\left\langle h^{2}\right\rangle  &=&\frac{1}{\widehat{q}\left( r\right) }%
\int_{0}^{2\pi }\frac{d\theta }{2\pi }\frac{r}{R_{0}}\frac{B_{\min }\left(
r\right) }{B_{P}}\left( \frac{B}{B_{\max }\left( r\right) }\right) ^{2} 
\nonumber \\
&=&\frac{1}{\widehat{q}\left( r\right) }\frac{B_{\min }\left( r\right) }{%
B_{\max }^{2}\left( r\right) }\int_{0}^{2\pi }\frac{d\theta }{2\pi }\frac{r}{%
R_{0}}\frac{B^{2}}{B_{P}}
\end{eqnarray}
and 
\begin{eqnarray}
\left\langle h^{2}\right\rangle  &=&\left( 1-\epsilon \right) ^{2}\frac{B_{P}%
}{BB_{0}}\int_{0}^{2\pi }\frac{d\theta }{2\pi }\frac{B^{2}}{B_{P}}  \nonumber
\\
&=&\left( 1-\epsilon \right) ^{2}\frac{1}{B_{0}}\int_{0}^{2\pi }\frac{%
d\theta }{2\pi }B  \nonumber \\
&=&\frac{\left( 1-\epsilon \right) ^{2}}{1+\epsilon }\int_{0}^{2\pi }\frac{%
d\theta }{2\pi }\Psi \left( r,\theta \right)   \nonumber \\
&=&\frac{\left( 1-\epsilon \right) ^{2}}{1+\epsilon }s^{*}  \nonumber \\
&=&\frac{\left( 1-\epsilon \right) ^{3/2}}{\left( 1+\epsilon \right) ^{1/2}}
\end{eqnarray}
with $s^{*}=\sqrt{\frac{1+\epsilon }{1-\epsilon }}$ as calculated in the
Appendix ****. Hence, in this simplified configuration, 
\begin{equation}
\left\langle \frac{1-\sqrt{1-h}\left( 1+\frac{h}{2}\right) }{h^{2}}%
\right\rangle =\left( 1+\epsilon \right) \int_{0}^{2\pi }\frac{1}{\Psi
\left( r,\theta \right) }\frac{1-\sqrt{1-\frac{\Psi \left( r,\theta \right) 
}{s^{*2}}}\left( 1+\frac{1}{2}\frac{\Psi \left( r,\theta \right) }{s^{*2}\,}%
\right) }{\left[ \frac{\Psi \left( r,\theta \right) }{s^{*2}\,}\right] ^{2}}%
\frac{d\theta }{2\pi }
\end{equation}
which becomes
\begin{equation}
\left\langle \frac{1-\sqrt{1-h}\left( 1+\frac{h}{2}\right) }{h^{2}}%
\right\rangle =\frac{\left( 1+\epsilon \right) ^{3}}{\left\langle 1-\epsilon
\right\rangle ^{2}}\int_{0}^{2\pi }\frac{1-\sqrt{1-\frac{\Psi \left(
r,\theta \right) }{s^{*2}}}\left( 1+\frac{1}{2}\frac{\Psi \left( r,\theta
\right) }{s^{*2}}\right) }{\Psi ^{3}\left( r,\theta \right) }\frac{d\theta }{%
2\pi }
\end{equation}

Since 
\begin{equation}
\left\{ 
\begin{array}{l}
\lim_{\epsilon \rightarrow 0}1-\left\langle h\right\rangle =\epsilon  \\ 
\lim_{\epsilon \rightarrow 0}1-\left\langle h^{2}\right\rangle =2\epsilon 
\end{array}
\right. 
\end{equation}
one obtains in a straightforward manner
\begin{equation}
\lim_{\epsilon \rightarrow 0}f_{tu,M}^{neo}\left( r\right) =\frac{3}{2}\sqrt{%
\epsilon }
\end{equation}
and 
\begin{equation}
\lim_{\epsilon \rightarrow 0}f_{tl,M}^{neo}\left( r\right) =\frac{3\sqrt{2}}{%
\pi }\sqrt{\epsilon }
\end{equation}

Using the interpolation formula (\ref{intertrapfrac}),

\begin{equation}
\lim_{\epsilon \rightarrow 0}\overline{f}_{t,M}^{neo}\left( r\right) =\left(
0.75\frac{3}{2}+0.25\frac{3\sqrt{2}}{\pi }\right) \sqrt{\epsilon }=1.46\sqrt{%
\epsilon }
\end{equation}

In this simplified configuration, it is interesting to evaluate directly the
thermal neoclassical expression, starting from (\ref{neotrapfracgenthermal}%
). One obtains
\begin{eqnarray}
\left\langle \sqrt{1-xh}\right\rangle  &=&\frac{1}{\widehat{q}\left(
r\right) }\int_{0}^{2\pi }\frac{d\theta }{2\pi }\epsilon \frac{B_{\min
}\left( r\right) }{B_{P}}\sqrt{1-x\frac{B}{B_{\max }}}  \nonumber \\
&=&\frac{\epsilon B}{\widehat{q}\left( r\right) B_{P}}\int_{0}^{2\pi }\frac{%
d\theta }{2\pi }\frac{B_{\min }\left( r\right) }{B}\sqrt{1-x\frac{B}{B_{\max
}}}  \nonumber \\
&=&\left( 1+\epsilon \right) \int_{0}^{2\pi }\frac{d\theta }{2\pi }\frac{%
B_{\min }\left( r\right) }{B}\sqrt{1-x\frac{B}{B_{\max }}}
\end{eqnarray}

Hence, 

\begin{equation}
\left\langle \sqrt{1-xh}\right\rangle =\int_{0}^{2\pi }\frac{d\theta }{2\pi }%
\left( 1+\epsilon \cos \theta \right) \sqrt{1-x\left( \frac{1-\epsilon }{%
1+\epsilon \cos \theta }\right) }
\end{equation}
or

\begin{equation}
\left\langle \sqrt{1-xh}\right\rangle =\int_{0}^{2\pi }\frac{d\theta }{2\pi }%
\sqrt{\left( 1+\epsilon \cos \theta \right) ^{2}-x\left( 1-\epsilon \right)
\left( 1+\epsilon \cos \theta \right) }
\end{equation}
and at the limit of large aspect ratio,
\begin{eqnarray}
\lim_{\epsilon \rightarrow 0}\left\langle \sqrt{1-xh}\right\rangle  &\simeq
&\int_{0}^{2\pi }\frac{d\theta }{2\pi }\sqrt{1+x\epsilon -x+(2-x)\epsilon
\cos \theta }  \nonumber \\
&=&\sqrt{1+x\epsilon -x}\int_{0}^{2\pi }\frac{d\theta }{2\pi }\sqrt{1+\frac{%
(2-x)\epsilon }{1+x\epsilon -x}\cos \theta }
\end{eqnarray}

If $\kappa =\frac{(2-x)\epsilon }{1+x\epsilon -x},$ with $0<\kappa <1$, one
must calculate the integral
\begin{eqnarray}
I &=&\int_{0}^{2\pi }\frac{d\theta }{2\pi }\sqrt{1+\kappa \cos \theta } 
\nonumber \\
&=&\int_{0}^{\pi }\frac{d\theta }{\pi }\sqrt{1+\kappa \cos 2\theta } 
\nonumber \\
&=&\int_{0}^{\pi }\frac{d\theta }{\pi }\sqrt{1+\kappa \left( 2\cos
^{2}\theta -1\right) }  \nonumber \\
&=&\int_{0}^{\pi }\frac{d\theta }{\pi }\sqrt{1+\kappa \left( 2\left( 1-\sin
^{2}\theta \right) -1\right) }  \nonumber \\
&=&\int_{0}^{\pi }\frac{d\theta }{\pi }\sqrt{1+\kappa -2\kappa \sin
^{2}\theta }  \nonumber \\
&=&\sqrt{1+\kappa }\int_{0}^{\pi }\frac{d\theta }{\pi }\sqrt{1-\frac{2\kappa 
}{1+\kappa }\sin ^{2}\theta }  \nonumber \\
&&\sqrt{1+\kappa }E\left( \pi \backslash \arcsin \left( \sqrt{\frac{2\kappa 
}{1+\kappa }}\right) \right) 
\end{eqnarray}
where $E\left( x\right) $ is the elliptic integral of second kind.

The trapped fraction $f_{t,M}^{neo}$ then becomes
\begin{equation}
\lim_{\epsilon \rightarrow 0}f_{t,M}^{neo}\left( r\right) \simeq 1-\frac{3}{4%
}\int_{0}^{1}\frac{x}{\sqrt{\frac{1+2\epsilon -x}{1+x\left( \epsilon
-1\right) }}E\left( \pi \backslash \arcsin \left( \sqrt{\frac{2\epsilon
\left( 2-x\right) }{1+2\epsilon -x}}\right) \right) }dx
\end{equation}

\subsection{Bremsstrahlung emission}

